\section{Limitations}
The strongest limitation is the participant count. The analysis contains 57
participants, including 19 dyslexia-labeled readers, so the result is a strong internal
confirmatory finding rather than an externally validated screening model. Labels are
operational research labels; they should not be interpreted as clinical diagnosis,
clinical screening, or medical validation. No independent external dataset is included,
and generalization beyond Danish natural reading remains open.

Text and speech exposure are not randomized, so the paper relies on exposure-only
comparisons, removal of exposure-count variables, and text-exposure audits. These
checks reduce the plausibility of a text-assignment explanation but do not replace an
independent balanced replication. BenchmarkBridge v1 adds internal EyeBench-style
unseen-reader and unseen-text checks, but exact official EyeBench folds were not
available in the local workspace; those results should therefore be described as
benchmark-relative rather than official benchmark scores. Calibration is reported, but
calibration estimates are limited by the small participant sample and should be treated
cautiously. Participant influence remains possible despite the leave-one-dyslexia-labeled
and remove-one-participant sensitivity checks.

LM alignment warnings are recorded and should be considered when interpreting
individual token-level features. Boundary-opacity labels are deterministic
orthographic proxies, not pronunciation-aware segmentation labels. Parser status is
\texttt{surface\_heuristic\_fallback}, so parser-syntax claims are deferred. Gemma
sensitivity is pending gated access and is not used in the main claim. Reviewer risks
are summarized in Table~\ref{tab:reviewer-risk}.

\begin{table}[t]
\centering
\small
\caption{Reviewer-risk summary and submission blocking status.}
\label{tab:reviewer-risk}
\begin{tabular}{llllll}
\toprule
risk & risk\_level & evidence\_available & remaining\_weakness & where\_to\_discuss & blocks\_submission \\
\midrule
small participant count & high & 57 participants; bootstrap/influence available & limited external power & main text and limitations & False \\
label provenance & high & label release documents operational labels & not clinical diagnosis & main text & False \\
text/speech exposure imbalance & medium & exposure-only DFM is weak; exposure audits included & not randomized text assignment & main text and appendix & False \\
reading-speed confound & medium & D3 excludes raw speed/global duration & other speed-like residuals possible & main text & False \\
DFM exposure vs sensitivity & low & D1 weak, D2/D3 strong & same dataset only & main text & False \\
leakage risk & low & LOPO and cross-fitted residualization & implementation complexity & appendix & False \\
calibration & medium & Brier and calibration slope recorded & small calibration sample & appendix & False \\
participant influence & medium & leave-one-dyslexia minimum ROC-AUC 0.8801 & small N remains & appendix & False \\
LM alignment warnings & medium & missingness/warnings recorded & DFM warning details need appendix & appendix & False \\
segmentation proxy limitations & medium & orthographic proxy only & not pronunciation-aware & limitations & False \\
parser fallback & medium & surface\_heuristic\_fallback documented & no true syntax & limitations & False \\
Gemma pending & medium & DFM locked; Gemma deferred & missing LM sensitivity & future work & False \\
\bottomrule
\end{tabular}
\end{table}

